\documentclass[11pt]{article}
\usepackage[margin=1in]{geometry}
\usepackage{amsmath,amssymb,amsfonts,mathtools,bm}
\usepackage[T1]{fontenc}
\usepackage[utf8]{inputenc}
\title{Inline and Display Mathematics}
\author{Source-backed grouped formula sample}
\date{}
\begin{document}
\maketitle
\section{Expressions}
\paragraph{Expression 1.} The following expression is taken from a source-backed formula corpus.

\begin{align*}\tilde{\gamma}^{a} p_{0a} = 0 = \tilde{\gamma}^{a} f^{\mu}{ }_ap_{0\mu} = (\tilde{\gamma}^{m} f^{\alpha}{ }_m +\tilde{\gamma}^{h} f^{\alpha}{ }_h )p_{0\alpha} + (\tilde{\gamma}^{m} f^{\sigma}{ }_m +\tilde{\gamma}^{h} f^{\sigma}{ }_h ) p_{0\sigma},\end{align*}

\paragraph{Expression 2.} The following expression is taken from a source-backed formula corpus.

\begin{align*}\psi_{-}(\vec{x}) = \frac{1}{\sqrt{2}} \frac{1}{i\partial_{-} -e V_{-}} \ast \xi (\vec{x}) - \frac{e}{\sqrt{2}}\frac{1}{\Delta_{\perp}} \ast (J_{pzm}^i (x_\perp) - q_{i} )\alpha^i \frac{1}{i\partial_{-} - e V_{-}} \ast\psi_{+} (\vec{x}) \;, \end{align*}

\paragraph{Expression 3.} The following expression is taken from a source-backed formula corpus.

\begin{align*}B_{\mu\nu}(t)=\sum_{\rho=0}^N\frac{\eta_{\mu\rho}\eta_{\nu\rho}}{\sqrt{8\omega_\mu}}\left[\left(1+\frac{\omega_\mu}{\Omega_\rho}\right)\rm{e}^{i\Omega_\rho t}+\left(1-\frac{\omega_\mu}{\Omega_\rho}\right)\rm{e}^{-i\Omega_\rho t}\right]\;.\end{align*}

\paragraph{Expression 4.} The following expression is taken from a source-backed formula corpus.

\begin{align*}\tilde{R}^{ab}=R^{ab}-K^a_c\wedge K^{cb}+\frac{1}{2}C^{ab}_{\;\cdot c}\tilde {T}^c+\frac{1}{4}\gamma^I_J \gamma^K_L\theta^a_I \theta^b_K A^J\wedge A^L+i\left(DK^{ab}-\frac{1}{4}\gamma^I_J(\theta^a_Ie^b+\theta^b_Ie^a)\wedge A^J\right). \end{align*}

\paragraph{Expression 5.} The following expression is taken from a source-backed formula corpus.

\begin{align*}\|\lambda(\phi)\|_{K,m+n}=\sup_{\stackrel{|p|\leq m+n}{z \in K}}|D^p (\lambda(\phi))(z)|\,\,,\qquad D^p=\frac{\partial^{|q|+|r|}}{\partial x^{q_1}_1 \cdots \partial x^{q_m}_m\partial\theta^{r_1}_1 \cdots \partial\theta^{r_n}_n}\,\,. \end{align*}
\end{document}
